\documentclass{situationist-preprint}

\title{The Van Meegeren Case}
\msid{AES-04-b}
\msfamily{AES-04}
\msdiscipline{Aesthetics}
\mstype{Companion manuscript}
\mskeywords{forgery; van Meegeren; Vermeer; artistic value; achievement; categories}

\begin{document}
\maketitle

\begin{abstract}
Han van Meegeren's forged Vermeers are the most cited case in the philosophy of forgery. AES-04-a argues that a perfect forgery is aesthetically identical to its original but lacks its artistic value, and that the difference in price tracks the latter. We examine how the case bears on that thesis. The forgeries were not perfect, and their faults became visible once they were seen in the right category; but the case still isolates achievement value, because the change in their standing in 1945 preceded any change in how they looked. The case thus confirms the account twice: once through the role of categories in aesthetic perception, and once through the dependence of artistic value on attributed achievement.
\end{abstract}

\section{The case}

In 1937 \emph{Christ and the Disciples at Emmaus} was acclaimed by the Vermeer scholar Abraham Bredius as an early masterpiece and was acquired for the Museum Boijmans in Rotterdam. It was painted by Han van Meegeren, who went on to produce further \enquote{Vermeers}, one of which was sold to Hermann G\"oring. Arrested in 1945 for having sold a national treasure to the enemy, van Meegeren declared that the paintings were his own, and to establish the claim he painted a further picture in the manner of Vermeer under supervision. He was tried in 1947 for forgery and fraud and convicted \citep{Lopez2008}. The paintings, which had been valued as among the most important Dutch pictures, lost almost all of their value overnight.

\section{Two lessons}

AES-04-a distinguishes aesthetic from artistic value (Definition~2.2) and argues that a perfect forgery (Definition~2.1) and its original have the same aesthetic value (Lemma~3.1) but differ in artistic value (Lemma~3.2), because only the original embodies the achievement in question (Principle~P2).

\emph{First lesson: categories.} The van Meegerens were not perfect forgeries. Viewers now commonly find the faces heavy-lidded and the drawing awkward, nothing like Vermeer, and remark on a resemblance to the film stills of the 1930s. This is what Principle~P3 of AES-04-a predicts. Aesthetic properties supervene on perceptible features as perceived in a category \citep{Walton1970}. Seen as Vermeers, and as members of a hypothesized early religious period of which no other examples were known, the pictures' peculiarities were absorbed as the marks of an immature manner. Seen as twentieth-century pastiches, the same peculiarities stand out. The change in perceived aesthetic quality is a change of category, not of eyesight.

\emph{Second lesson: achievement.} The collapse in value did not wait for the change in perception; it followed the confession. In 1945 nothing about the paintings' appearance had changed, and many viewers had not yet learned to see them differently. What changed was the achievement they could be credited with. A picture that showed Vermeer working out a religious manner of which there was no other record embodied, if genuine, an achievement of great art-historical significance; the same picture as van Meegeren's showed the imitation of a manner already known. This is the pattern Proposition~4.1 of AES-04-a attributes to prices: the premium follows attributions of achievement, and here it fell with them before the aesthetic reappraisal began.

\begin{objection}[The forger's achievement]
Van Meegeren deceived the leading experts of his day. That is an achievement too, and on Principle~P2 it should confer artistic value on the paintings.
\end{objection}

\begin{reply}
It is an achievement, but not the one for which the paintings were valued, and not an artistic achievement of the relevant kind. Deceiving experts is a success in imitation and in exploiting expectations about an artist's development; the value it confers is that of a historical curiosity, and the paintings now command the modest prices, and the documentary interest, that fit it. \citet{Dutton1979} put the point by saying that a forgery misrepresents the achievement it embodies; the van Meegerens embody a real achievement, but a different one from what they were taken to embody, and their value changed accordingly.
\end{reply}

\section{Conclusion}

The van Meegeren case does not test the claim that perfect forgeries are aesthetically identical to their originals, since the forgeries were not perfect. It confirms the two other parts of AES-04-a: that aesthetic qualities are perceived in categories, which explains how the forgeries passed and why they no longer do, and that artistic value and price follow achievement, which explains why their value fell before anyone saw them differently.

\begin{thebibliography}{99}\small
\bibitem[Dutton(1979)]{Dutton1979} Dutton, D. (1979). Artistic crimes: The problem of forgery in the arts. \emph{The British Journal of Aesthetics}, 19(4), 302--314.
\bibitem[Lopez(2008)]{Lopez2008} Lopez, J. (2008). \emph{The Man Who Made Vermeers: Unvarnishing the Legend of Master Forger Han van Meegeren}. Harcourt.
\bibitem[Walton(1970)]{Walton1970} Walton, K.~L. (1970). Categories of art. \emph{The Philosophical Review}, 79(3), 334--367.
\end{thebibliography}

\end{document}
